%%! Author = jdubec
%! Date = 16/02/2026

\documentclass[11pt, a4paper]{article}

% ─── Packages ────────────────────────────────────────────────────────────────
\usepackage[utf8]{inputenc}
\usepackage[T1]{fontenc}
\usepackage[margin=2.5cm]{geometry}
\usepackage{graphicx}
\usepackage{hyperref}
\usepackage{booktabs}       % nicer tables (\toprule, \midrule, \bottomrule)
\usepackage{float}          % [H] placement for figures and tables
\usepackage{xcolor}         % colors (used by listings)
\usepackage{listings}       % code listings
\hypersetup{colorlinks=true, linkcolor=black, citecolor=black, urlcolor=blue!70!black}

% ─── Code listing style ─────────────────────────────────────────────────────
\lstset{
    basicstyle=\ttfamily\small,
    breaklines=true,
    frame=single,
    numbers=left,
    numberstyle=\tiny\color{gray},
    xleftmargin=2em,
    framexleftmargin=1.5em
}
% Pro tip: for syntax highlighting, look into the 'minted' package.

% ─── Bibliography ────────────────────────────────────────────────────────────
\usepackage[style=numeric, backend=bibtex]{biblatex}
\addbibresource{references.bib}

\setcounter{biburlnumpenalty}{9000}
\setcounter{biburlucpenalty}{9000}
\setcounter{biburllcpenalty}{9000}

% ─── Language (swap to lang/sk for Slovak) ───────────────────────────────────
% Slovak language pack
\usepackage[slovak]{babel}

\newcommand{\LabelTitle}{Návrh databázového systému}
\newcommand{\LabelSession}{Cvičenie}

\newcommand{\LabelCourse}{Databázové technológie}
\newcommand{\LabelCode}{DBS}
\newcommand{\LabelUniversity}{Slovenská technická univerzita v~Bratislave}
\newcommand{\LabelFaculty}{Fakulta informatiky a~informačných technológií}
\newcommand{\LabelAssignment}{Zadanie 4: Letisko}
\newcommand{\LabelStudent}{Študent}
\newcommand{\LabelGroup}{Študijná skupina}
\newcommand{\LabelDate}{Dátum}
\newcommand{\LabelYear}{Akademický rok}
\newcommand{\LabelContents}{Obsah}

% ─── Student details ─────────────────────────────────────────────────────────
\newcommand{\StudentName}{Filip Múdry, Dániel Ruzsik}
\newcommand{\TeacherName}{Ing. Lukáš Rucek}
\newcommand{\StudyGroup}{Skupina 8}
\newcommand{\StudySession}{Streda 08:00, \TeacherName}
\newcommand{\AcademicYear}{2025/2026}

% =============================================================================
\begin{document}

% ─── Cover page ──────────────────────────────────────────────────────────────
\begin{titlepage}
    \centering

    \includegraphics[width=0.3\textwidth]{assets/logo.png}

    \vspace{0.3cm}
    {\large \LabelUniversity}\\[2pt]
    {\large \LabelFaculty}

    \vspace{2.5cm}
    {\Huge\bfseries \LabelTitle}

    \vspace{0.6cm}
    {\LARGE \LabelCode{} --- \LabelCourse}

    \vspace{0.4cm}
    {\Large \LabelAssignment}

    \vfill

    \begin{tabular}{r l}
        \textbf{\LabelStudent:} & \StudentName  \\[4pt]
        \textbf{\LabelGroup:}   & \StudyGroup   \\[4pt]
        \textbf{\LabelSession:}      &
        \StudySession \\[4pt]
        \textbf{\LabelYear:}    & \AcademicYear \\[4pt]
        \textbf{\LabelDate:}    & \today
        
        
        \\
    \end{tabular}

    \vspace{1.5cm}
\end{titlepage}

% ─── Table of contents ───────────────────────────────────────────────────────
\renewcommand{\contentsname}{\LabelContents}
\tableofcontents
\newpage

% =============================================================================
% This document serves as both the assignment template and a quick guide
% to the LaTeX features you will need. Replace the content, keep the structure.
%
% Template source: https://github.com/FIIT-Databases/assignment-template
% =============================================================================

\section{Špecifikácia zadania}
\label{sec:specifikacia}

Toto zadanie je zamerané na modelovanie prostredia letiska tak, aby databáza
neodrážala iba technické fungovanie infraštruktúry, ale predovšetkým reálnu
skúsenosť cestujúceho počas celej cesty. Cieľom projektu nie je detailná
simulácia interných letiskových technológií, ale návrh dátového modelu, ktorý
umožní prehľadne sledovať cestujúceho od odchodu na letisko až po úspešný
nástup do lietadla.

V praxi to znamená, že sa sústredíme na celý používateľský tok: cestujúci si
kupuje letenku, má priradené miesto v lietadle, môže mať batožinu, presúva sa
cez konkrétne letisko a je naviazaný na konkrétny let. Zároveň chceme vedieť,
na akej trase sa daný let nachádza, či ide o priamy let alebo prestup, aká je
zdrojová a cieľová destinácia a aký je aktuálny stav lietadla, ktoré let
obsluhuje.

Dôležitou súčasťou zadania je aj orientácia na aktuálnu polohu a stav lietadla.
Používateľ alebo operátor systému má vedieť zistiť, či sa lietadlo nachádza v
hangári, je priradené na runway, alebo je asociované s aktívnym letom. Tento
pohľad je dôležitý najmä pri prestupoch a plánovaní nadväzných spojení, kde
časová a stavová informácia rozhoduje o tom, či cestujúci stihne ďalší úsek
cesty.


Výsledkom má byť databázový model, ktorý bude dostatočne konzistentný na
implementáciu informačného systému pre sledovanie cesty cestujúceho, plánovanie
letov a evidenciu vzťahov medzi letiskami, lietadlami, letmi a cestujúcimi.

\subsection*{Otázka 1}
\textbf{Aký problém systém rieši a pre koho je určený? Aká je jeho pridaná hodnota?}

Systém rieši evidenciu základných informácií o fungovaní letiska. Umožňuje
ukladať a spravovať údaje o cestujúcich, letoch, letenkách, lietadlách a
batožine. Cieľom je mať všetky tieto informácie prepojené tak, aby bolo možné
jednoducho zistiť napríklad na aký let je cestujúci priradený, aké lietadlo
daný let vykonáva alebo akú batožinu má cestujúci evidovanú.

Systém je určený hlavne pre pracovníkov letiska, ktorí potrebujú pracovať s
údajmi o letoch a cestujúcich. Môžu to byť napríklad pracovníci pri odbavení
cestujúcich, pracovníci správy letov alebo administrátori systému.

Pridaná hodnota systému je v tom, že všetky dôležité údaje sú uložené na jednom
mieste a navzájom prepojené. Vďaka tomu je možné rýchlo získať potrebné
informácie a znižuje sa riziko chýb pri práci s dátami.

\subsection*{Otázka 2}
\textbf{Kto sú aktéri systému? Aké majú role, zodpovednosti a úrovne prístupu?}

Aktérmi systému sú cestujúci, piloti, pracovníci letiska a administrátor systému.

Cestujúci je hlavný subjekt evidovaný v systéme. Systém uchováva jeho osobné
údaje a informácie o letenkách, batožine a priradenom lete. Celý návrh databázy
je z veľkej časti postavený okolo evidencie cestujúceho a jeho pohybu v rámci
letiska až po nástup do lietadla.

Pilot je aktér, ktorý je priradený ku konkrétnemu letu. Systém eviduje základné
informácie o pilotovi a umožňuje priradiť pilota k letu, ktorý vykonáva. Pilot
je teda zodpovedný za samotnú realizáciu letu.

Pracovníci letiska pracujú so systémom pri každodennej prevádzke. Môžu evidovať
údaje o letoch, lietadlách, runway, hangároch alebo spravovať informácie o
cestujúcich, letenkách a batožine. Ich prístup je obmedzený len na tie časti
systému, ktoré potrebujú pre svoju prácu.

Administrátor systému má najvyššiu úroveň prístupu. Zodpovedá za správu
databázy, kontrolu správnosti údajov a technickú údržbu systému. Má prístup ku
všetkým údajom v databáze a zabezpečuje, aby systém fungoval správne.

\subsection*{Otázka 3}
\textbf{Aké sú kľúčové entity, ktoré systém musí evidovať? Prečo sú dôležité?}

Kľúčové entity systému predstavujú hlavné objekty, ktoré sú potrebné na
evidenciu letov a pohybu cestujúcich v rámci letiska.
\newline

Jednou z najdôležitejších entít je \textbf{passenger}, ktorá reprezentuje
cestujúceho. Táto entita uchováva základné údaje o osobe, ktorá využíva let,
napríklad meno, email alebo číslo pasu.

Ďalšou kľúčovou entitou je \textbf{ticket}, ktorá predstavuje letenku
cestujúceho. Tá prepája cestujúceho s konkrétnym letom a obsahuje informácie
ako cena, číslo sedadla alebo stav letenky.

Entita \textbf{flight} reprezentuje samotný let. Uchováva informácie o čase
odletu a príletu, priradenom lietadle, pilotovi a runway. Táto entita je
dôležitá, pretože prepája viacero ďalších entít v systéme.

Entita \textbf{airplane} reprezentuje lietadlo, ktoré vykonáva let. Uchováva
informácie napríklad o leteckej spoločnosti, modeli lietadla alebo jeho
kapacite.

Entita \textbf{airport} reprezentuje letisko, z ktorého let odlieta alebo na
ktoré prichádza. Uchováva základné informácie ako názov letiska, kód, mesto a
krajinu.

Entita \textbf{route} reprezentuje trasu letu medzi dvoma letiskami. Umožňuje
evidovať, z ktorého letiska let odlieta a na ktoré smeruje.

Entita \textbf{baggage} eviduje batožinu cestujúceho, ktorá je priradená ku
konkrétnej letenke.
\newline

Tieto entity sú kľúčové, pretože spolu reprezentujú hlavný proces systému --
evidenciu cestujúceho, jeho letenky, letu a súvisiacich informácií.

\subsection*{Otázka 4}
\textbf{Aké biznis pravidlá a obmedzenia musí systém rešpektovať? Ako sa premietnu do integritných obmedzení?}

Systém musí rešpektovať viaceré biznis pravidlá, ktoré zabezpečujú správnosť a
logiku údajov uložených v databáze. Tieto pravidlá sa v databáze implementujú
pomocou rôznych integritných obmedzení, napríklad kontrolných podmienok
(\texttt{CHECK}), unikátnych obmedzení alebo vzťahov medzi tabuľkami.
\newline
Jedným z pravidiel je, že cena letenky nemôže byť záporná. Letenka musí mať
vždy kladnú alebo nulovú hodnotu. Toto pravidlo je možné implementovať pomocou
kontrolného obmedzenia napríklad:

\begin{lstlisting}[language=SQL]
CHECK (price >= 0)
\end{lstlisting}

Ďalším pravidlom je, že čas príletu musí byť vždy neskôr ako čas odletu. Let
nemôže skončiť skôr, ako začal. Toto pravidlo sa dá zabezpečiť pomocou
kontrolného obmedzenia:

\begin{lstlisting}[language=SQL]
CHECK (arrival_time > departure_time)
\end{lstlisting}

Systém tiež zabezpečuje, že hmotnosť batožiny musí byť väčšia ako nula. Nie je
možné evidovať batožinu s nulovou alebo zápornou hmotnosťou. Toto pravidlo je
možné implementovať napríklad pomocou:

\begin{lstlisting}[language=SQL]
CHECK (weight > 0)
\end{lstlisting}

Ďalším pravidlom je, že kapacita lietadla musí byť kladné číslo, pretože
lietadlo musí mať aspoň jedno sedadlo. Toto sa dá implementovať napríklad
pomocou:

\begin{lstlisting}[language=SQL]
CHECK (capacity > 0)
\end{lstlisting}

Systém tiež zabezpečuje, že dve osoby nemôžu mať rovnaké sedadlo na tom istom
lete. Jedno sedadlo môže byť obsadené len jedným cestujúcim. Toto pravidlo sa
dá implementovať pomocou unikátneho obmedzenia napríklad na kombináciu:

\begin{lstlisting}[language=SQL]
UNIQUE (flight_id, seat_id)
\end{lstlisting}

Ďalším pravidlom je, že zdrojové a cieľové letisko v trase letu nemôžu byť
rovnaké, pretože let musí smerovať z jedného letiska na iné. Toto pravidlo je
možné implementovať napríklad pomocou:

\begin{lstlisting}[language=SQL]
CHECK (src_airport_id <> dest_airport_id)
\end{lstlisting}

Tieto biznis pravidlá zabezpečujú, že systém nebude umožňovať ukladať
nelogické alebo neplatné údaje a že údaje v databáze budú konzistentné a
realistické vzhľadom na fungovanie letiska.


\section{Definovanie používateľských scenárov}
\label{sec:usecase}

Táto sekcia obsahuje päť realistických používateľských scenárov, ktoré
reprezentujú typické operácie systému. Scenáre sú navrhnuté tak, aby neboli
iba jednoduché CRUD operácie, ale pokrývali celý proces od interakcie aktéra až
po zmeny na dátovej vrstve.

\subsection{UC1 - Rezervácia letu a vystavenie letenky}

\subsubsection*{Stručný popis}

Scenár pokrýva proces, v ktorom cestujúci vyhľadá vhodný let, dokončí rezerváciu
a systém mu vystaví letenku so sedadlom.

\subsubsection*{Hlavný tok}
\begin{enumerate}
    \item Cestujúci zadá odletové a cieľové letisko, dátum a preferencie.
    \item Systém zobrazí dostupné lety.
    \item Cestujúci vyberie konkrétny let.
    \item Systém zobrazí dostupné sedadlá a cenu.
    \item Cestujúci vyberie sedadlo.
    \item Systém uloží voľbu sedadla.
    \item Cestujúci potvrdí rezerváciu a vykoná platbu.
    \item Systém  vytvorí letenku.
    \item Cestujúci dostane potvrdenie rezervácie.
\end{enumerate}

\subsubsection*{Čo sa deje na dátovej vrstve}
\begin{itemize}
    \item \textbf{Aké tabuľky sa prehľadávajú:} systém najskôr pracuje s tabuľkami \texttt{airport}, \texttt{route} a \texttt{flight}, aby našiel vhodné spojenia podľa odletu, cieľa a času. Následne sa číta \texttt{airplane} kvôli kapacite a \texttt{ticket} kvôli už obsadeným sedadlám. V kroku potvrdenia sa používa aj \texttt{passenger}, aby sa nová rezervácia korektne naviazala na konkrétneho cestujúceho.
    \item \textbf{Aké entity vznikajú:} po úspešnom potvrdení rezervácie vzniká nová doménová entita letenky, teda nový záznam v \texttt{ticket}. Tento záznam reprezentuje výsledok celého procesu rezervácie a obsahuje väzbu na let aj cestujúceho. Podľa implementácie môže popri tom vzniknúť aj technický záznam do auditnej tabuľky (napríklad kto, kedy a pre aký let vytvoril rezerváciu).
    \item \textbf{Aké agregácie sa vykonávajú:} systém priebežne počíta, koľko miest je už predaných alebo rezervovaných pre konkrétny let, a tento počet porovnáva s kapacitou lietadla. Zároveň môže dopočítať počet voľných miest, ktorý sa zobrazuje používateľovi pri výbere sedadla. Tieto agregácie sú dôležité, aby sa zabránilo prebookovaniu a aby používateľ videl len reálne dostupné možnosti.
    \item \textbf{Aké záznamy sa menia:} mení sa najmä stav novej letenky (napr. rezervovaná, zaplatená), priradenie sedadla, čas vytvorenia rezervácie a prípadne ďalšie atribúty súvisiace s platbou. Ak je systém navrhnutý s dočasnou rezerváciou pred platbou, môže sa počas procesu meniť aj interný stav tej istej letenky pred jej finálnym potvrdením.
    \item \textbf{Kontroly konzistencie:} overuje sa, že vybrané sedadlo je stále voľné pre daný let a že rezervácia spĺňa integritné pravidlá. Kľúčová je kontrola \texttt{UNIQUE (flight\_id, seat\_id)} a tiež validácia cien cez \texttt{CHECK (price >= 0)}. Rovnako sa kontrolujú cudzie kľúče, aby \texttt{ticket} vždy odkazoval na existujúci \texttt{passenger} a existujúci \texttt{flight}.
\end{itemize}

\subsubsection*{Alternatívne toky}
\begin{itemize}
    \item \textbf{Ak nie je splnená podmienka voľných miest:} let je vypredaný, systém neuloží rezerváciu a ponúkne iný let alebo termín.
    \item \textbf{Ak nie je splnená podmienka dostupnosti sedadla:} sedadlo bolo medzičasom obsadené, systém vráti používateľa na výber sedadla.
    \item \textbf{Ak nie je splnená podmienka úspešnej platby:} platba zlyhá, záznam letenky sa nevytvorí alebo zostane v neaktívnom stave a rezervácia sa nedokončí.
\end{itemize}

\subsubsection*{Podporný UML diagram}

\textbf{Sekvenčný diagram} je vhodný, pretože zachytáva komunikáciu medzi
cestujúcim, rezervačným modulom, platobným modulom a databázou v presnom
časovom poradí.

\IfFileExists{assets/UC1.jpg}{%
\begin{figure}[H]
    \centering
    \includegraphics[width=0.75\textwidth]{assets/UC1.jpg}
    \caption{Podporný UML diagram pre UC1}
\end{figure}
}{%
\textit{Obrázok assets/UC1.jpg nebol nájdený.}
}

\subsection{UC2 - Zobrazenie trasy letu a cieľovej destinácie}

\subsubsection*{Stručný popis}

Scenár pokrýva situáciu, keď cestujúci alebo pracovník letiska potrebuje získať
prehľad o trase letu, prestupoch a finálnej destinácii.

\subsubsection*{Hlavný tok}
\begin{enumerate}
    \item Pracovník otvorí detail letu alebo letenky.
    \item Systém načíta údaje o trase letu.
    \item Systém zobrazí odletové a príletové letisko.
    \item Pracovník vidí kompletný itinerár vrátane cieľovej destinácie.
\end{enumerate}

\subsubsection*{Čo sa deje na dátovej vrstve}
\begin{itemize}
    \item \textbf{Aké tabuľky sa prehľadávajú:} hlavným vstupom je \texttt{ticket} alebo identifikátor letu, odkiaľ sa pokračuje na \texttt{flight}. Na základe letu sa číta \texttt{route}, ktorá určuje zdrojové a cieľové letisko, a následne sa z \texttt{airport} načítajú názvy, kódy a ďalšie detaily letísk pre zobrazenie používateľovi.
    \item \textbf{Aké entity vznikajú:} v štandardnom priebehu nevzniká nová biznis entita, pretože ide o čítací scenár. Systém však môže vytvoriť technický záznam o tom, že bol detail itinerára zobrazený (audit, log prístupov, telemetry event), čo pomáha pri diagnostike alebo bezpečnostnom dohľade.
    \item \textbf{Aké agregácie sa vykonávajú:} ak trasa obsahuje viac segmentov, systém ich zoradí do správneho poradia a vypočíta súhrnné údaje, napríklad počet prestupov alebo celkovú dĺžku cesty. Pri priamom lete je agregácia jednoduchšia, no logika zostáva rovnaká: pripraviť konzistentný výstup pre UI v jednom prehľade.
    \item \textbf{Aké záznamy sa menia:} doménové tabuľky ako \texttt{ticket}, \texttt{flight}, \texttt{route} a \texttt{airport} sa nemenia, keďže scenár je read-only. Ak je zapnutý audit, mení sa iba technická vrstva logovania (nový záznam o zobrazení alebo aktualizácia počítadla prístupov).
    \item \textbf{Kontroly konzistencie:} systém overuje, že väzby \texttt{ticket -> flight -> route} existujú a že všetky referencované letiská sú v tabuľke \texttt{airport}. Pri viacsegmentovej trase sa navyše kontroluje časová a logická nadväznosť úsekov, aby nevznikol chybný itinerár s nelogickým poradím alebo rozpadnutými väzbami.
\end{itemize}

\subsubsection*{Alternatívne toky}
\begin{itemize}
    \item Trasa nie je dostupná: systém zobrazí stav „trasa nie je definovaná".
\end{itemize}

\subsubsection*{Podporný UML diagram}

\textbf{Diagram aktivít} je vhodný, pretože prehľadne zobrazuje vetvenie
(priamy let vs. prestup) a rozhodovacie body pri skladaní itinerára.

\IfFileExists{assets/UC2.jpg}{%
\begin{figure}[H]
    \centering
    \includegraphics[width=0.95\textwidth]{assets/UC2.jpg}
    \caption{Podporný UML diagram pre UC2}
\end{figure}
}{%
\textit{Obrázok assets/UC2.jpg nebol nájdený.}
}

\subsection{UC3 - Sledovanie aktuálneho stavu letu a lietadla}

\subsubsection*{Stručný popis}

Scenár pokrýva monitorovanie aktuálneho stavu letu a prevádzkového stavu
priradeného lietadla (napr. hangár, runway, aktívny let).

\subsubsection*{Hlavný tok}
\begin{enumerate}
    \item Používateľ vyberie let, ktorý chce sledovať.
    \item Systém načíta priradené lietadlo.
    \item Systém zistí aktuálny stav letu (napr. boarding, delayed, departed).
    \item Systém vyhodnotí, či je lietadlo v hangári, na runway alebo v aktívnom
    lete.
    \item Systém zobrazí konsolidovaný status letu aj lietadla.
\end{enumerate}

\subsubsection*{Čo sa deje na dátovej vrstve}
\begin{itemize}
    \item \textbf{Aké tabuľky sa prehľadávajú:} jadrom scenára sú \texttt{flight} a \texttt{airplane}, ktoré určujú, aký let sa sleduje a aké lietadlo je k nemu aktuálne priradené. Doplňujúco sa čítajú \texttt{runway}, \texttt{airplane\_hangar} a \texttt{hangar}, aby bolo možné určiť prevádzkový kontext lietadla. Ak je dostupná tabuľka stavových udalostí, systém z nej berie najnovšie zmeny statusu.
    \item \textbf{Aké entity vznikajú:} vo väčšine prípadov nevznikajú nové doménové entity, pretože ide o monitorovací scenár. Môže však vzniknúť technický záznam monitorovania, napríklad snapshot aktuálneho stavu alebo záznam o tom, že operátor otvoril živý status konkrétneho letu.
    \item \textbf{Aké agregácie sa vykonávajú:} systém agreguje viacero informačných zdrojov do jedného výsledného „live statusu“ — kombinuje stav letu, stav lietadla a jeho umiestnenie v infraštruktúre. Prakticky ide o zostavenie jednotného prevádzkového pohľadu, ktorý je použiteľný pre dispečing aj informačné panely.
    \item \textbf{Aké záznamy sa menia:} pri príchode nových prevádzkových dát sa aktualizujú stavy v \texttt{flight} a/alebo \texttt{airplane}, prípadne sa pridáva nový stavový event. Zmeny sú typicky časovo citlivé, preto je dôležité udržiavať posledný známy validný status a jeho časovú značku.
    \item \textbf{Kontroly konzistencie:} kontroluje sa, že väzby medzi letom, lietadlom a infraštruktúrou sú platné a že status nie je v rozpore s fyzickým stavom (napr. lietadlo nemôže byť súčasne v hangári aj v aktívnom odlete). Zároveň sa overuje čerstvosť údajov, aby sa neprezentovali zastarané informácie ako aktuálny stav.
\end{itemize}

\subsubsection*{Alternatívne toky}
\begin{itemize}
    \item Lietadlo nie je priradené: systém vráti informáciu o chýbajúcom
    priradení.
    \item Stavové údaje sú neaktuálne: systém zobrazí posledný dostupný stav s
    časovou značkou.
\end{itemize}

\subsubsection*{Podporný UML diagram}

\textbf{Sekvenčný diagram} je vhodný, pretože ukazuje postupné dopyty na viaceré
prevádzkové entity a následné zlúčenie výsledkov do jedného statusu.

\IfFileExists{assets/UC3.jpg}{%
\begin{figure}[H]
    \centering
    \includegraphics[width=0.95\textwidth]{assets/UC3.jpg}
    \caption{Podporný UML diagram pre UC3}
\end{figure}
}{%
\textit{Obrázok assets/UC3.jpg nebol nájdený.}
}

\subsection{UC4 - Odbavenie cestujúceho a evidovanie batožiny}

\subsubsection*{Stručný popis}

Scenár pokrýva check-in proces, v ktorom sa overí letenka cestujúceho, zaznamená
sa odbavenie a eviduje sa odovzdaná batožina.

\subsubsection*{Hlavný tok}
\begin{enumerate}
    \item Cestujúci sa identifikuje na check-ine.
    \item Systém overí platnosť letenky a priradenie k letu.
    \item Pracovník zaznamená batožinu a jej hmotnosť.
    \item Systém uloží záznam o batožine.
    \item Systém nastaví stav letenky cestujúceho na odbavený (checked in).
\end{enumerate}

\subsubsection*{Čo sa deje na dátovej vrstve}
\begin{itemize}
    \item \textbf{Aké tabuľky sa prehľadávajú:} systém číta \texttt{passenger} a \texttt{ticket}, aby overil identitu a platnosť letenky, následne \texttt{flight}, aby zistil, či je check-in pre daný let otvorený. Pri evidencii batožiny sa používa \texttt{baggage}, kde sa kontrolujú existujúce kusy a pripravuje zápis nových údajov.
    \item \textbf{Aké entity vznikajú:} pri úspešnom odbavení vzniká jeden alebo viac nových záznamov v \texttt{baggage}, podľa počtu odovzdaných kusov batožiny. Tieto entity sú naviazané na letenku/cestujúceho a predstavujú oficiálnu evidenciu batožiny pre konkrétny let.
    \item \textbf{Aké agregácie sa vykonávajú:} systém počíta celkovú hmotnosť batožiny na cestujúceho (alebo na letenku) a porovnáva ju s limitom povoleným tarifou či pravidlami pre daný let. Pri viacerých kusoch sa agreguje súčet hmotností a podľa výsledku sa rozhoduje o doplatku alebo potrebe manuálneho zásahu.
    \item \textbf{Aké záznamy sa menia:} mení sa stav letenky na checked-in, dopĺňa sa čas odbavenia a podľa návrhu aj identifikátor pracovníka, ktorý check-in vykonal. Popri tom sa ukladajú nové záznamy batožiny a môžu sa meniť pomocné polia súvisiace s odbavením (napr. interné poznámky alebo stav spracovania).
    \item \textbf{Kontroly konzistencie:} overuje sa validita letenky, príslušnosť k letu, otvorené check-in okno a pravidlá kvality údajov pre batožinu, najmä \texttt{CHECK (weight > 0)}. Ak hmotnosť presiahne limit, systém musí vyvolať alternatívny tok (doplatok, korekcia údajov alebo zamietnutie ďalšieho kusu).
\end{itemize}

\subsubsection*{Alternatívne toky}
\begin{itemize}
    \item Neplatná letenka: systém check-in zamietne.
    \item Nadlimitná batožina: systém vyžiada doplatok alebo úpravu batožiny.
    \item Let je uzavretý pre check-in: systém odbavenie neumožní.
\end{itemize}

\subsubsection*{Podporný UML diagram}

\textbf{Diagram aktivít} je vhodný, pretože obsahuje viacero rozhodnutí
(validita letenky, hmotnosť batožiny, otvorenosť check-inu).

\IfFileExists{assets/UC4.jpg}{%
\begin{figure}[H]
    \centering
    \includegraphics[width=0.95\textwidth]{assets/UC4.jpg}
    \caption{Podporný UML diagram pre UC4}
\end{figure}
}{%
\textit{Obrázok assets/UC4.jpg nebol nájdený.}
}

\subsection{UC5 - Boarding cestujúceho do lietadla}

\subsubsection*{Stručný popis}

Scenár pokrýva nástup cestujúceho do lietadla pri bráne po úspešnom odbavení a
bezpečnostnej kontrole.

\subsubsection*{Hlavný tok}
\begin{enumerate}
    \item Cestujúci sa dostaví k bráne a predloží boarding pass.
    \item Systém overí, že letenka patrí k správnemu letu a je v stave „checked-in".
    \item Systém overí, že boarding je pre daný let otvorený.
    \item Pracovník potvrdí nástup cestujúceho.
    \item Systém zaznamená cestujúceho ako „boarded".
\end{enumerate}

\subsubsection*{Čo sa deje na dátovej vrstve}
\begin{itemize}
    \item \textbf{Aké tabuľky sa prehľadávajú:} pri načítaní boarding passu sa číta \texttt{ticket} (stav letenky), \texttt{flight} (stav a časové okno boardingu) a \texttt{passenger} (zhoda identity). Ak systém vedie samostatnú tabuľku boarding udalostí, tá sa používa na kontrolu predchádzajúcich vstupov a na uloženie nového výsledku.
    \item \textbf{Aké entity vznikajú:} vzniká záznam boarding udalosti s časom skenu, výsledkom validácie a prípadným dôvodom zamietnutia. Táto entita tvorí auditnú stopu nástupu a je dôležitá pre spätné dohľadanie priebehu boardingu.
    \item \textbf{Aké agregácie sa vykonávajú:} systém priebežne počíta počet nastúpených cestujúcich na daný let a porovnáva ho s počtom checked-in alebo predaných leteniek. Agregácia umožňuje operatívne sledovať priebeh nástupu a včas identifikovať nezrovnalosti pred uzatvorením brány.
    \item \textbf{Aké záznamy sa menia:} po úspešnom boardingu sa letenka aktualizuje na stav boarded a doplní sa čas nástupu. Súčasne sa zapíše boarding event; pri neúspešnom pokuse sa môže uložiť zamietnutý event bez zmeny stavu letenky.
    \item \textbf{Kontroly konzistencie:} kontroluje sa zhoda boarding passu s cestujúcim a letom, otvorené boarding okno a idempotencia (ten istý cestujúci nesmie byť zapísaný ako boarded viackrát). Tieto kontroly zabraňujú duplicitám, chybnej evidencii a neoprávnenému nástupu.
\end{itemize}

\subsubsection*{Alternatívne toky}
\begin{itemize}
    \item Boarding pass je neplatný alebo pre iný let: nástup sa zamietne.
    \item Cestujúci nie je odbavený: systém vráti cestujúceho na check-in.
    \item Brána je uzavretá: nástup nie je povolený.
\end{itemize}

\subsubsection*{Podporný UML diagram}

\textbf{Sekvenčný diagram} je vhodný, pretože presne zachytáva interakciu medzi
pracovníkom brány, boarding systémom a databázou počas krátkej, ale kritickej
operácie.

\IfFileExists{assets/UC5.jpg}{%
\begin{figure}[H]
    \centering
    \includegraphics[width=0.95\textwidth]{assets/UC5.jpg}
    \caption{Podporný UML diagram pre UC5}
\end{figure}
}{%
\textit{Obrázok assets/UC5.jpg nebol nájdený.}
}

\section{Návrh relačného modelu}
\label{sec:relacnymodel}

\subsection{Diagram relačného modelu}

Relačný model je zdokumentovaný aj graficky. Diagram je vložený priamo v tomto
PDF a zároveň je dostupný ako samostatný obrázok vo formáte JPG.

\IfFileExists{assets/Fyzicky_datovy_model.jpg}{%
\begin{figure}[H]
    \centering
    \includegraphics[width=0.98\textwidth]{assets/Fyzicky_datovy_model.jpg}
    \caption{Fyzický diagram relačného modelu}
    \label{fig:fyzicky-datovy-model}
\end{figure}
}{%
\textit{Obrázok assets/Fyzicky\_datovy\_model.jpg nebol nájdený.}
}

\subsection{Detailný popis každej tabuľky}

\subsubsection*{\texttt{airport}}
\begin{itemize}
    \item \textbf{Účel:} evidencia letísk (názov, kód, mesto, krajina).
    \item \textbf{Atribúty a obmedzenia:} \texttt{id SERIAL PK}, \texttt{name VARCHAR}, \texttt{code
        VARCHAR}, \texttt{city VARCHAR}, \texttt{country VARCHAR}.
    \item \textbf{Primárny kľúč:} umelý kľúč \texttt{id}; je stabilný a nezávislý od biznis hodnoty
        \texttt{code}.
    \item \textbf{Cudzie kľúče:} tabuľka nemá vlastné FK; odkazujú na ňu \texttt{runway} a \texttt{route}.
\end{itemize}

\subsubsection*{\texttt{passenger}}
\begin{itemize}
    \item \textbf{Účel:} evidencia cestujúcich pre rezervácie, check-in a boarding.
    \item \textbf{Atribúty a obmedzenia:} \texttt{id SERIAL PK}, \texttt{firstname VARCHAR},
        \texttt{lastname VARCHAR}, \texttt{email VARCHAR}, \texttt{passport\_number VARCHAR},
        \texttt{nationality VARCHAR}.
    \item \textbf{Primárny kľúč:} umelý kľúč \texttt{id}; praktickejší a bezpečnejší ako použitie čísla pasu.
    \item \textbf{Cudzie kľúče:} tabuľka nemá vlastné FK.
\end{itemize}

\subsubsection*{\texttt{ticket}}
\begin{itemize}
    \item \textbf{Účel:} reprezentuje letenku/rezerváciu cestujúceho na konkrétny let.
    \item \textbf{Atribúty a obmedzenia:} \texttt{id SERIAL PK}, \texttt{flight\_id INT FK},
        \texttt{passenger\_id INT FK}, \texttt{seat\_id INT FK}, \texttt{ticket\_status\_id INT FK},
        \texttt{price DECIMAL CHECK (price >= 0)}, \texttt{booked\_at TIMESTAMPTZ DEFAULT
        CURRENT\_TIMESTAMP}, \texttt{checked\_in\_at TIMESTAMPTZ}.
    \item \textbf{Primárny kľúč:} umelý kľúč \texttt{id}; letenka je samostatná transakčná entita.
    \item \textbf{Cudzie kľúče:} \texttt{flight\_id -> flight(id)}, \texttt{passenger\_id ->
        passenger(id)}, \texttt{seat\_id -> seat(id)}, \texttt{ticket\_status\_id ->
        ticket\_status(id)}.
\end{itemize}

\subsubsection*{\texttt{ticket\_status}}
\begin{itemize}
    \item \textbf{Účel:} číselník stavov leteniek.
    \item \textbf{Atribúty a obmedzenia:} \texttt{id SERIAL PK}, \texttt{status VARCHAR}.
    \item \textbf{Primárny kľúč:} umelý kľúč \texttt{id}; štandardné riešenie pre číselníky.
    \item \textbf{Cudzie kľúče:} tabuľka nemá vlastné FK.
\end{itemize}

\subsubsection*{\texttt{flight}}
\begin{itemize}
    \item \textbf{Účel:} evidencia konkrétnych letov a ich prevádzkových väzieb.
    \item \textbf{Atribúty a obmedzenia:} \texttt{id SERIAL PK}, \texttt{runway\_id INT FK},
        \texttt{airplane\_id INT FK}, \texttt{pilot\_id INT FK}, \texttt{flight\_status\_id INT FK},
        \texttt{flight\_number VARCHAR}, \texttt{departure\_at TIMESTAMPTZ}, \texttt{arrival\_at
        TIMESTAMPTZ}.
    \item \textbf{Primárny kľúč:} umelý kľúč \texttt{id}; číslo letu nemusí byť globálne jedinečné v čase.
    \item \textbf{Cudzie kľúče:} \texttt{runway\_id -> runway(id)}, \texttt{airplane\_id ->
        airplane(id)}, \texttt{pilot\_id -> pilot(id)}, \texttt{flight\_status\_id ->
        flight\_status(id)}.
\end{itemize}

\subsubsection*{\texttt{flight\_status}}
\begin{itemize}
    \item \textbf{Účel:} číselník stavov letu (napr. departed, landed, cancelled).
    \item \textbf{Atribúty a obmedzenia:} \texttt{id SERIAL PK}, \texttt{status VARCHAR}.
    \item \textbf{Primárny kľúč:} umelý kľúč \texttt{id}.
    \item \textbf{Cudzie kľúče:} tabuľka nemá vlastné FK.
\end{itemize}

\subsubsection*{\texttt{airplane}}
\begin{itemize}
    \item \textbf{Účel:} evidencia lietadiel použitých na lety.
    \item \textbf{Atribúty a obmedzenia:} \texttt{id SERIAL PK}, \texttt{airline VARCHAR},
        \texttt{model VARCHAR}, \texttt{capacity INT CHECK (capacity > 0)}, \texttt{status VARCHAR
        DEFAULT 'active'}.
    \item \textbf{Primárny kľúč:} umelý kľúč \texttt{id}; model nevyužíva prirodzený jednoznačný
        identifikátor.
    \item \textbf{Cudzie kľúče:} tabuľka nemá vlastné FK.
\end{itemize}

\subsubsection*{\texttt{baggage}}
\begin{itemize}
    \item \textbf{Účel:} evidencia batožiny priradenej ku konkrétnej letenke.
    \item \textbf{Atribúty a obmedzenia:} \texttt{id SERIAL PK}, \texttt{ticket\_id INT FK},
        \texttt{baggage\_status\_id INT FK}, \texttt{weight REAL CHECK (weight > 0)}.
    \item \textbf{Primárny kľúč:} umelý kľúč \texttt{id}; umožňuje evidovať viac kusov k jednej letenke.
    \item \textbf{Cudzie kľúče:} \texttt{ticket\_id -> ticket(id)}, \texttt{baggage\_status\_id ->
        baggage\_status(id)}.
\end{itemize}

\subsubsection*{\texttt{baggage\_status}}
\begin{itemize}
    \item \textbf{Účel:} číselník stavov batožiny (napr. naložená, stratená).
    \item \textbf{Atribúty a obmedzenia:} \texttt{id SERIAL PK}, \texttt{status VARCHAR}.
    \item \textbf{Primárny kľúč:} umelý kľúč \texttt{id}.
    \item \textbf{Cudzie kľúče:} tabuľka nemá vlastné FK.
\end{itemize}

\subsubsection*{\texttt{runway}}
\begin{itemize}
    \item \textbf{Účel:} evidencia dráh prislúchajúcich letiskám.
    \item \textbf{Atribúty a obmedzenia:} \texttt{id SERIAL PK}, \texttt{airport\_id INT FK},
        \texttt{name VARCHAR}, \texttt{status VARCHAR DEFAULT 'active'}.
    \item \textbf{Primárny kľúč:} umelý kľúč \texttt{id}.
    \item \textbf{Cudzie kľúče:} \texttt{airport\_id -> airport(id)}.
\end{itemize}

\subsubsection*{\texttt{airplane\_hangar}}
\begin{itemize}
    \item \textbf{Účel:} väzobná tabuľka M:N medzi lietadlami a hangármi s históriou pobytu.
    \item \textbf{Atribúty a obmedzenia:} \texttt{id SERIAL PK}, \texttt{airplane\_id INT FK},
        \texttt{hangar\_id INT FK}, \texttt{arrived\_at TIMESTAMPTZ DEFAULT CURRENT\_TIMESTAMP},
        \texttt{departed\_at TIMESTAMP}.
    \item \textbf{Primárny kľúč:} uprednostnený umelý kľúč \texttt{id}; umožňuje opakované
        historické pobyty rovnakého lietadla v tom istom hangári.
    \item \textbf{Cudzie kľúče:} \texttt{airplane\_id -> airplane(id)}, \texttt{hangar\_id -> hangar(id)}.
\end{itemize}

\subsubsection*{\texttt{hangar}}
\begin{itemize}
    \item \textbf{Účel:} evidencia hangárov na letiskách.
    \item \textbf{Atribúty a obmedzenia:} \texttt{id SERIAL PK}, \texttt{airport\_id INT FK},
        \texttt{name VARCHAR}, \texttt{capacity INT CHECK (capacity > 0)}, \texttt{status VARCHAR
        DEFAULT 'operational'}.
    \item \textbf{Primárny kľúč:} umelý kľúč \texttt{id}.
    \item \textbf{Cudzie kľúče:} \texttt{airport\_id -> airport(id)}.
\end{itemize}

\subsubsection*{\texttt{route}}
\begin{itemize}
    \item \textbf{Účel:} evidencia trasy letu medzi zdrojovým a cieľovým letiskom.
    \item \textbf{Atribúty a obmedzenia:} \texttt{id SERIAL PK}, \texttt{dest\_airport\_id INT FK},
        \texttt{src\_airport\_id INT FK}, \texttt{flight\_id INT FK}.
    \item \textbf{Primárny kľúč:} umelý kľúč \texttt{id}; jednoduchá identifikácia každého segmentu trasy.
    \item \textbf{Cudzie kľúče:} \texttt{flight\_id -> flight(id)}, \texttt{src\_airport\_id ->
        airport(id)}, \texttt{dest\_airport\_id -> airport(id)}.
\end{itemize}

\subsubsection*{\texttt{boarding}}
\begin{itemize}
    \item \textbf{Účel:} evidencia nástupu cestujúceho na palubu.
    \item \textbf{Atribúty a obmedzenia:} \texttt{id SERIAL PK}, \texttt{ticket\_id INT FK},
        \texttt{boarded\_at TIMESTAMPTZ DEFAULT CURRENT\_TIMESTAMP}, \texttt{boarding\_status\_id
        INT FK}.
    \item \textbf{Primárny kľúč:} umelý kľúč \texttt{id}.
    \item \textbf{Cudzie kľúče:} \texttt{ticket\_id -> ticket(id)}, \texttt{boarding\_status\_id ->
        boarding\_status(id)}.
\end{itemize}

\subsubsection*{\texttt{boarding\_status}}
\begin{itemize}
    \item \textbf{Účel:} číselník stavov boardingu.
    \item \textbf{Atribúty a obmedzenia:} \texttt{id INT/SERIAL PK}, \texttt{status VARCHAR}.
    \item \textbf{Primárny kľúč:} umelý kľúč \texttt{id}.
    \item \textbf{Cudzie kľúče:} tabuľka nemá vlastné FK.
\end{itemize}

\subsubsection*{\texttt{pilot}}
\begin{itemize}
    \item \textbf{Účel:} evidencia pilotov priraditeľných k letom.
    \item \textbf{Atribúty a obmedzenia:} \texttt{id SERIAL PK}, \texttt{firstname VARCHAR},
        \texttt{lastname VARCHAR}, \texttt{licence\_number VARCHAR}.
    \item \textbf{Primárny kľúč:} umelý kľúč \texttt{id}; praktickejší ako použitie
        \texttt{licence\_number} ako PK.
    \item \textbf{Cudzie kľúče:} tabuľka nemá vlastné FK.
\end{itemize}

\subsubsection*{\texttt{seat}}
\begin{itemize}
    \item \textbf{Účel:} evidencia sedadiel v jednotlivých lietadlách.
    \item \textbf{Atribúty a obmedzenia:} \texttt{id SERIAL PK}, \texttt{seat\_number VARCHAR},
        \texttt{airplane\_id INT FK}.
    \item \textbf{Primárny kľúč:} umelý kľúč \texttt{id}; prirodzený kandidát je kombinácia
        \texttt{(airplane\_id, seat\_number)}.
    \item \textbf{Cudzie kľúče:} \texttt{airplane\_id -> airplane(id)}.
\end{itemize}

\subsection{Popis vzťahov medzi tabuľkami}

Väčšina vzťahov v tomto modeli je typu 1:N a na fyzickej úrovni je
realizovaná cudzím kľúčom v podradenej tabuľke. Jediný vzťah typu M:N je
medzi \texttt{airplane} a \texttt{hangar}, ktorý je realizovaný cez väzobnú
tabuľku \texttt{airplane\_hangar}.

\subsubsection*{Vzťah medzi \texttt{passenger} a \texttt{ticket} (1:N)}
Jeden pasažier môže mať viac leteniek, ale každá letenka patrí práve jednému
pasažierovi. Tento vzťah je realizovaný cez \texttt{ticket.passenger\_id}.

\subsubsection*{Vzťah medzi \texttt{flight} a \texttt{ticket} (1:N)}
Jeden let môže mať viac leteniek, pričom každá letenka patrí jednému letu.
Realizácia je cez \texttt{ticket.flight\_id}.

\subsubsection*{Vzťah medzi \texttt{seat} a \texttt{ticket} (1:N)}
Podľa diagramu môže byť jedno sedadlo použité vo viacerých letenkách v čase,
ale každá letenka odkazuje na jedno sedadlo. Tento vzťah je realizovaný cez
\texttt{ticket.seat\_id}.

\subsubsection*{Vzťah medzi \texttt{ticket\_status} a \texttt{ticket} (1:N)}
Jeden stav môže patriť viacerým letenkám, ale každá letenka má práve jeden
stav. Vzťah je realizovaný cez \texttt{ticket.ticket\_status\_id}.

\subsubsection*{Vzťah medzi \texttt{ticket} a \texttt{baggage} (1:N)}
Jedna letenka môže mať viac kusov batožiny, ale každý kus batožiny patrí jednej
letenke. Realizácia je cez \texttt{baggage.ticket\_id}.

\subsubsection*{Vzťah medzi \texttt{baggage\_status} a \texttt{baggage} (1:N)}
Jeden stav batožiny môže byť priradený viacerým kusom batožiny, pričom každý
kus batožiny má jeden stav. Realizácia je cez
\texttt{baggage.baggage\_status\_id}.

\subsubsection*{Vzťah medzi \texttt{ticket} a \texttt{boarding} (1:N)}
Podľa diagramu môže mať jedna letenka viac boarding záznamov a každý boarding
záznam patrí jednej letenke. Tento vzťah je realizovaný cez
\texttt{boarding.ticket\_id}. Z vecného hľadiska by sa možno očakával skôr
jeden boarding záznam na jednu letenku, ale diagram to takto striktne
neobmedzuje.

\subsubsection*{Vzťah medzi \texttt{boarding\_status} a \texttt{boarding} (1:N)}
Jeden stav boardingu môže patriť viacerým boarding záznamom, ale každý boarding
záznam má práve jeden stav. Realizácia je cez
\texttt{boarding.boarding\_status\_id}.

\subsubsection*{Vzťah medzi \texttt{airport} a \texttt{runway} (1:N)}
Jedno letisko môže mať viac dráh, ale každá dráha patrí jednému letisku.
Realizácia je cez \texttt{runway.airport\_id}.

\subsubsection*{Vzťah medzi \texttt{runway} a \texttt{flight} (1:N)}
Jedna dráha môže byť použitá pre viac letov, ale každý let má pridelenú jednu
dráhu. Realizácia je cez \texttt{flight.runway\_id}.

\subsubsection*{Vzťah medzi \texttt{airplane} a \texttt{flight} (1:N)}
Jedno lietadlo môže vykonať viac letov, ale každý let má pridelené jedno
lietadlo. Realizácia je cez \texttt{flight.airplane\_id}.

\subsubsection*{Vzťah medzi \texttt{pilot} a \texttt{flight} (1:N)}
Jeden pilot môže pilotovať viac letov, ale každý let má jedného pilota.
Realizácia je cez \texttt{flight.pilot\_id}.

\subsubsection*{Vzťah medzi \texttt{flight\_status} a \texttt{flight} (1:N)}
Jeden stav letu môže byť priradený viacerým letom, ale každý let má presne
jeden stav. Realizácia je cez \texttt{flight.flight\_status\_id}.

\subsubsection*{Vzťah medzi \texttt{airport} a \texttt{hangar} (1:N)}
Jedno letisko môže mať viac hangárov, ale každý hangár patrí jednému letisku.
Realizácia je cez \texttt{hangar.airport\_id}.

\subsubsection*{Vzťah medzi \texttt{airplane} a \texttt{seat} (1:N)}
Jedno lietadlo má viac sedadiel, ale každé sedadlo patrí jednému lietadlu.
Realizácia je cez \texttt{seat.airplane\_id}.

\subsubsection*{Vzťah medzi \texttt{flight} a \texttt{route} (1:N)}
Podľa diagramu môže mať jeden let viac trás a každá trasa patrí jednému letu.
Je to trochu netypické, pretože bežne by sa očakávalo, že let má práve jednu
trasu. Na fyzickej úrovni je vzťah realizovaný cez \texttt{route.flight\_id}.

\subsubsection*{Vzťah medzi \texttt{airport} a \texttt{route} (1:N v dvoch roliach)}
Existujú dva samostatné vzťahy 1:N: jeden v role zdrojového letiska cez
\texttt{route.src\_airport\_id} a druhý v role cieľového letiska cez
\texttt{route.dest\_airport\_id}.

\subsubsection*{Vzťah medzi \texttt{airplane} a \texttt{hangar} (M:N)}
Vzťah M:N je realizovaný cez tabuľku \texttt{airplane\_hangar}. Jedno lietadlo
môže byť v rôznych hangároch a jeden hangár môže obsahovať rôzne lietadlá.
Navyše táto väzobná tabuľka uchováva aj vlastné atribúty
\texttt{arrived\_at} a \texttt{departed\_at}, takže ide o typickú väzbu s
vlastnými údajmi.

\subsection{Kardinalita, povinnosť a správanie pri mazaní}

\begin{itemize}
    \item V modeli sa typicky vyskytuje kardinalita \texttt{1} na strane rodiča a \texttt{0..*} na
        strane potomka.
    \item Z pohľadu zachovania histórie je vo väčšine doménových vzťahov bezpečnejšie použiť
        \texttt{RESTRICT} (alebo ekvivalent \texttt{NO ACTION}) pri mazaní rodiča.
    \item \texttt{CASCADE} sa dá zvážiť len pri čisto závislých pomocných záznamoch (napr.
        \texttt{ticket -> baggage}, \texttt{ticket -> boarding}) a len ak systém nechce dlhodobo
        držať históriu.
    \item Pri transakčných a auditných údajoch (letenky, lety, boarding, batožina) je odporúčaný
        konzervatívny prístup s \texttt{RESTRICT}.
\end{itemize}

\subsection{Ako sú zabezpečené integritné obmedzenia z kapitoly 2}

\begin{itemize}
    \item \textbf{Entitná integrita:} každá tabuľka má primárny kľúč, takže každý záznam je
        jednoznačne identifikovateľný.
    \item \textbf{Referenčná integrita:} cudzími kľúčmi je vynútená existencia rodičovských záznamov
        (napr. \texttt{ticket} bez platného \texttt{passenger}/\texttt{flight} nie je možné uložiť).
    \item \textbf{Doménová integrita:} dátové typy, \texttt{CHECK} podmienky a \texttt{DEFAULT}
        hodnoty (\texttt{price >= 0}, \texttt{weight > 0}, \texttt{capacity > 0},
        \texttt{CURRENT\_TIMESTAMP}, predvolené stavy).
    \item \textbf{Dôležité doplnenia:} odporúča sa doplniť \texttt{UNIQUE} obmedzenia pre
        \texttt{airport.code}, \texttt{passenger.email}, \texttt{passenger.passport\_number},
        \texttt{pilot.licence\_number} a \texttt{seat(airplane\_id, seat\_number)}.
    \item \textbf{Ďalšie odporúčané pravidlá:} \texttt{route.src\_airport\_id <>
        route.dest\_airport\_id}, \texttt{flight.departure\_at < flight.arrival\_at} a doménové
        pravidlo zhody sedadla s lietadlom letu (typicky trigger alebo rozšírené obmedzenie).
\end{itemize}

\subsection{Je možné realizovať všetky používateľské scenáre?}

Vo všeobecnosti \textbf{áno}. Navrhnuté tabuľky a väzby pokrývajú kľúčové scenáre:
rezerváciu letenky, správu letu a trasy, check-in, boarding, evidenciu batožiny,
a tiež prevádzkový kontext lietadiel (dráhy, hangáre).

Zároveň platí, že pre prísnejšie biznis pravidlá je vhodné model doplniť o
ďalšie databázové obmedzenia. Bez nich databáza sama negarantuje napríklad
jednoznačný predaj sedadla na konkrétny let, presný počet povolených boarding
záznamov na letenku alebo striktne jednu trasu pre každý let (ak je to
požiadavka domény).


\subsection{Zhodnotenie}


\subsubsection*{Ktorá časť riešenia nám zabrala najviac času a prečo?}

Čo sa týka návrhu celkovej štruktúry zadanie, pretože najprv sme zvažovali celkový návrh z hľadiska infraštruktúry letiska ale postupom času sme zmenili túto možnosť na návrh z pohľadu bežného cestujúceho teda od samotného príchodu na letisko až ku odletu, prestupmi letov a hlavne evidenciou jednotlivých pasažierov, lietadiel, pilotov, batožiny a iného. Jednoducho sme chceli mať prehľad na všetkým nie len sa pasažier check-inom a hotovo. Takže to je dôvod, prečo nám to zobralo najviac času, lebo sme to prerábali na iný návrh. 

\subsubsection*{S čím sme si nevedeli poradiť a ako sme to vyriešili?}

\textbf{Prestupy pri ceste s viacerými letiskami}

Pri návrhu databázového modelu sme narazili na problém, ako reprezentovať cestu cestujúceho s prestupom. 
V reálnom systéme môže cestujúci absolvovať viac letov na jednej ceste. Napríklad, ak cestujúci cestuje z Bratislavy do New Yorku, cesta môže vyzerať nasledovne:
\begin{center}
    Bratislava → Frankfurt → New York. \\
\end{center} 

Najprv sme uvažovali o použití jedného \texttt{ticket}-u na viac letov. Tento prístup by však vyžadoval vytvorenie dodatočnej väzobnej tabuľky medzi \texttt{ticket} a \texttt{flight}. Ďalšia možnosť by mohla byť rozšírenie tabuľky \texttt{route} redundantnými stĺpcami (zhodujúcimi so stĺpcami tabuľky \texttt{flight}).

Nakoniec sme sa rozhodli pre jednoduchšie riešenie, ktoré používa \textbf{jeden ticket} na reprezentovanie \textbf{iba jedenéhu letu}. V prípade prestupu cestujúci nakúpi viac leteniek, pričom každá je priradená ku konkrétnemu letu. Našli sme, že tento prístup bude jednoduchší na implementáciu a zároveň postačuje pre potreby nášho systému.

Nie je to vec, s ktorou by sme si nevedeli rady, ale dumali sme nad tým, či chceme mať perfektne simulované "real-life" riešenie letiska, alebo si trochu uľahčíme prácu, nakoľko už máme dostatočný počet potrebných tabuliek na absolvovanie zadania.

\subsubsection*{Čo bolo pre nás úplne nové?}

Oproti minulým rokom sa nám páčilo prepojenie návrhu aj s potrebnými prípadmi použitia, pretože sme si vedeli náš návrh skontrolovať pri reálnom produkčnom chode takéhoto navrhovaného systému. 

\subsubsection*{Pokrýva naše riešenie správne zvolenú tému? Čo sme zjednodušili alebo vynechali?}

Riešenie pokrýva zvolenú tému správne, pretože zahŕňa kľúčové procesy
cestujúceho od rezervácie až po nástup do lietadla a zároveň modeluje
prevádzkové súvislosti letu. 
Zjednodušenie nastalo už v spomínanom odstránení, alebo vynechaní tabuľky \texttt{JOURNEY}, ktorá bola vykonaná z dôvodu lepšej evidencie prechodov pri jednotlivých letoch. Z hľadiska komplexnosti riešenia a dosiahnutia potrebného počtu zmysluplných tabuliek sme túto možnosť zamietli.

\subsubsection*{Aký bol náš pracovný proces? Postupovali sme iteratívne?}

Ako prvé, sme si zadefinovali doménu systému a kľúčové entity. 
Tým sme si vybudovali svet, v rámci čoho bude naše riešenie existovať.

Následne sme vytvorili prvú verziu väzieb a pravidiel, ktoré sme neskôr postupne upravovali na základe používateľských scenárov. 
Po každej iterácii sme kontrolovali, či návrh stále podporuje všetky kľúčové operácie systému bez porušenia konzistencia dát.

\subsubsection*{Použili sme AI nástroje? Ak áno – na čo presne, aký bol výstup a ako sme ho upravovali?}

Áno, AI nástroje sme použili ako podporný prostriedok na štruktúrovanie textu,
spresnenie formulácií a návrh alternatívnych verzií popisov use-case scenárov. Takisto aj na kontrolu správnosti modelu, ktorý sme sami navrhli, aby sme predišli problémom, ktoré mohli nastať v neskoršej časti zadania. Výstupy sme vždy manuálne kontrolovali, upravovali podľa zadania a
prispôsobovali konkrétnemu dátovému modelu, aby výsledný text zodpovedal
reálnemu návrhu a neobsahoval všeobecné alebo nepresné tvrdenia.

\subsubsection*{Prečo sme si vybrali túto rámcovú tému?}

Tému letiskového systému sme si vybrali paradoxne preto, lebo ani jeden z nás nikdy neletel lietadlom. Chceli sme mať lepší prehľad o tom, ako fungujú letiská a procesy pri nástupe do lietadla. 

\subsection{Zdroje}

Zdroje budú doplnené po finalizácii zadania a overení doménových pravidiel.


% \printbibliography









\section{Analysis}
\label{sec:analysis}

The goal of this project was to create a submission template for the DBS
course that students can clone and start writing in immediately. The template
should be easy to use, support Slovak and English, and not require any prior
\LaTeX{} experience\footnote{If this is your first time with \LaTeX{}, the
Overleaf guide \cite{overleaf-docs} covers everything you need in about 30 minutes.}.

The template repository is available at \cite{template-repo}.

% ── Subsections ──────────────────────────────────────────────────────────────
% Use \subsection{} and \subsubsection{} when you need more structure.

\subsection{Requirements}

We identified the following requirements:

\begin{itemize}
    \item Clean cover page with faculty branding.
    \item Bilingual support (SK/EN) via a single toggle.
    \item Examples of tables, figures, code, and citations.
    \item Minimal package dependencies --- it should compile everywhere.
\end{itemize}

When listing ordered steps, use \texttt{enumerate} instead:

\begin{enumerate}
    \item Clone the repository.
    \item Edit \texttt{main.tex}.
    \item Run \texttt{make}.
    \item Submit the PDF.
\end{enumerate}


\newpage
\section{Proposal}
\label{sec:proposal}

We compared a few approaches before settling on the current structure.
Table~\ref{tab:comparison} summarizes our options.

% ── Tables ───────────────────────────────────────────────────────────────────
% Use \begin{table}[H] to place the table exactly here.
% Label it with \label{} and reference it with Table~\ref{}.

\begin{table}[H]
    \centering
    \caption{Comparison of template approaches}
    \label{tab:comparison}
    \begin{tabular}{l c c c}
        \toprule
        \textbf{Approach}  & \textbf{Easy to edit} & \textbf{Bilingual} & \textbf{Compiles anywhere} \\
        \midrule
        Word document       & yes                   & manual             & n/a  \\
        Markdown + Pandoc   & yes                   & no                 & maybe \\
        \LaTeX{} template   & yes                   & yes                & yes  \\
        \bottomrule
    \end{tabular}
\end{table}

We chose \LaTeX{} because it handles structured documents, citations,
and cross-references well --- things that matter when writing technical
reports\footnote{It also looks way better than a Word document. But we are
biased.}. Unrelated: the wood frog can survive being frozen solid for months ---
its heart stops, it stops breathing, and in spring it just thaws and hops away.
We found this less stressful than debugging \LaTeX{} table alignment.


\newpage
\section{Implementation}
\label{sec:implementation}

The template lives in a simple directory structure with localization files
separated from the main document. Switching language requires changing one
line:

% ── Code listings ────────────────────────────────────────────────────────────
% Use lstlisting for short code snippets.

\begin{lstlisting}[language=TeX, caption={Switching language in the preamble}]
% English
% English language pack
\usepackage[english]{babel}

\newcommand{\LabelTitle}{Database Model Design}
\newcommand{\LabelCourse}{Database Technologies}
\newcommand{\LabelSession}{Exercise}

\newcommand{\LabelCode}{DBS}
\newcommand{\LabelUniversity}{Slovak University of Technology in Bratislava}
\newcommand{\LabelFaculty}{Faculty of Informatics and Information Technologies}
\newcommand{\LabelAssignment}{Assignment 4: Airport}
\newcommand{\LabelStudent}{Student}
\newcommand{\LabelGroup}{Study group}
\newcommand{\LabelDate}{Date}
\newcommand{\LabelYear}{Academic year}
\newcommand{\LabelContents}{Contents}

% Slovak
% Slovak language pack
\usepackage[slovak]{babel}

\newcommand{\LabelTitle}{Návrh databázového systému}
\newcommand{\LabelSession}{Cvičenie}

\newcommand{\LabelCourse}{Databázové technológie}
\newcommand{\LabelCode}{DBS}
\newcommand{\LabelUniversity}{Slovenská technická univerzita v~Bratislave}
\newcommand{\LabelFaculty}{Fakulta informatiky a~informačných technológií}
\newcommand{\LabelAssignment}{Zadanie 4: Letisko}
\newcommand{\LabelStudent}{Študent}
\newcommand{\LabelGroup}{Študijná skupina}
\newcommand{\LabelDate}{Dátum}
\newcommand{\LabelYear}{Akademický rok}
\newcommand{\LabelContents}{Obsah}
\end{lstlisting}

For SQL or other languages, change the \texttt{language} parameter:

\begin{lstlisting}[language=SQL, caption={Example SQL query}]
SELECT s.name, COUNT(*) AS submission_count
FROM students s
JOIN submissions sub ON s.id = sub.student_id
GROUP BY s.name
ORDER BY submission_count DESC;
\end{lstlisting}

% ── Inline code ──────────────────────────────────────────────────────────────
% Use \texttt{} for inline code: \texttt{SELECT * FROM ...}

References to tables, figures, and sections are automatic. For example,
this is Section~\ref{sec:implementation}, the comparison is in
Table~\ref{tab:comparison}, and the faculty logo is shown in
Figure~\ref{fig:logo} in Attachment~\ref{sec:attachment-a1}.
If you add or reorder content, the numbers update on the next build.

% ── Citations ────────────────────────────────────────────────────────────────
% Add entries to references.bib, then cite them with \cite{key}.
% The bibliography at the end is generated automatically.

The database fundamentals referenced throughout the course are covered
in \cite{elmasri2016}. For PostgreSQL-specific syntax, see
\cite{postgresql-docs}.


\newpage
\section{Challenges}
\label{sec:challenges}

The main challenge was keeping the template simple enough that a student
seeing \LaTeX{} for the first time would not immediately close the file.
We solved this by limiting the preamble to essential packages and
putting all examples in one place --- this document.

We also learned that BibTeX interprets \texttt{@} signs in comments as
entry types, which caused confusing errors. Small things like this are
easy to fix once you know about them, but hard to debug the first time\footnote{A
group of frogs is called an ``army.'' A group of students debugging BibTeX at
2\,AM is called ``Tuesday.''}.


\newpage
\section{Conclusion}
\label{sec:conclusion}

The template meets the requirements from Section~\ref{sec:analysis}.
It compiles with a standard \LaTeX{} distribution, supports two
languages, and provides working examples of every feature students
are likely to need.

For anything not covered here, the Overleaf documentation
\cite{overleaf-docs} is an excellent reference.

% ─── Bibliography ────────────────────────────────────────────────────────────
\printbibliography

% ─── Attachments ─────────────────────────────────────────────────────────────
\newpage
\appendix

\section{Sample figures}
\label{sec:attachment-a1}

% ── Figures ──────────────────────────────────────────────────────────────────
% Use \begin{figure}[H] to place the figure exactly here.
% Label it with \label{} and reference it with Figure~\ref{}.
% Keep images in the assets/ directory.

\begin{figure}[H]
    \centering
    \includegraphics[width=0.25\textwidth]{assets/logo.png}
    \caption{Faculty logo used on the cover page}
    \label{fig:logo}
\end{figure}

Place your diagrams, screenshots, and other supporting images in this
section. Reference them from the main text using
\texttt{Figure\~{}\textbackslash ref\{fig:logo\}}.

Some frogs can absorb water through a patch of skin on their belly called a
``drinking patch'' --- they literally sit in puddles to hydrate. This has
nothing to do with databases but is objectively more interesting than
anything on this page.

% You can control the image size with the width parameter:
%   width=0.5\textwidth   ... half the page width
%   width=5cm             ... exact size
%   scale=0.8             ... 80% of original

% To add more figures, copy the block above and change the file path,
% caption, and label.

\end{document}